\subsection{4. Đối tượng và phạm vi nghiên cứu}

\indent\indent \textbf{Đối tượng nghiên cứu:}
\begin{itemize}
    \item Các dấu vết làm giả văn bản trên ảnh tài liệu (vết cắt dán viền, vân bề mặt, sai lệch phông chữ, sai lệch căn dòng, đứt gãy tính hòa nhập).
    \item Kỹ thuật trích xuất đặc trưng không gian bằng mạng tích chập gọn nhẹ EfficientNet-B0.
    \item Kỹ thuật phân tích đặc trưng miền tần số cục bộ bằng biến đổi BlockDCT $8 \times 8$.
    \item Cơ chế học đa nhiệm (Multi-task Learning) kết hợp phân đoạn mặt nạ pixel và phân loại nhãn.
\end{itemize}
\vspace{0.3cm}

\textbf{Phạm vi nghiên cứu:}
\begin{itemize}
    \item \textbf{Về dữ liệu:} Sử dụng tập dữ liệu Minh-TFR (được trích xuất từ DTD-lightweight và TFR) gồm các mẫu ảnh tài liệu kèm mặt nạ ground-truth và nhãn phân loại 6 đặc trưng.
    \item \textbf{Về mô hình:} Tập trung vào kiến trúc học sâu LightweightDTD với tổng tham số dưới 5 triệu tham số, tối ưu hóa cho tốc độ suy luận nhanh trên GPU/CPU.
    \item \textbf{Về thực nghiệm:} Cài đặt trên môi trường PyTorch, đánh giá trên tập kiểm thử độc lập 2,000 mẫu (`Minh-TFR TestSet`).
\end{itemize}
